Estos tres términos se confunden fácilmente porque suenan parecido, pero cada uno impone una restricción distinta. La diferencia está en dos preguntas: ¿tiene que ser contiguo? y ¿importa el orden?

\begin{center}
\begin{tabular}{|l|c|c|p{6.5cm}|}
\hline
\textbf{Término} & \textbf{¿Contiguo?} & \textbf{¿Orden importa?} & \textbf{Ejemplo con \texttt{[1,2,3,4]}} \\
\hline
Subarreglo & Sí & Sí & \texttt{[2,3]} es válido. \texttt{[1,3]} \textbf{no} lo es (salta el 2, deja de ser contiguo). \\
\hline
Subsecuencia & No & Sí (se mantiene el orden relativo) & \texttt{[1,3]} es válido (se salta el 2). \texttt{[3,1]} \textbf{no} lo es (invierte el orden). \\
\hline
Subconjunto & No & No importa & \texttt{\{1,3\}} y \texttt{\{3,1\}} son el mismo subconjunto. \\
\hline
\end{tabular}
\end{center}

En strings los mismos conceptos aplican con otro nombre: \textbf{substring} (o \textit{subcadena}) es el equivalente a subarreglo (contiguo), y \textbf{subsecuencia de un string} es el equivalente a subsecuencia (no contiguo, orden preservado). Ejemplo con \texttt{"hello"}: \texttt{"ell"} es substring; \texttt{"hlo"} no es substring pero sí es subsecuencia (h, l, o aparecen en ese orden, saltando letras).

\textit{Regla práctica:} si el problema dice "contiguo" o habla de un "rango \([l,r]\)", es subarreglo/substring. Si dice "elige algunos elementos sin cambiar su orden", es subsecuencia. Si dice "elige un subconjunto" o el orden de elección no se menciona como relevante, es subconjunto (se suele recorrer con máscaras de bits, \(2^n\) posibilidades).
